% =====================================================================
% 08_prop4.tex — Proposition 4: equilibrium portfolios (WP-D)
% =====================================================================
\section{Proposition 4: equilibrium portfolios}
\label{sec:prop4}

This section proves Proposition~4 of KL: the comparative statics of the
equilibrium international portfolio --- Home's holdings of capital (a claim to the
world equity return) and its issuance of safe dollar bonds --- with respect to
the supply of safe dollar debt $b^g_{H,s}$ and to the cross-country risk-tolerance
ratio $\rraF/\rra$. The proof combines the international-portfolio
equation~\eqref{eq:eq99} of Section~\ref{sec:portfolio-engine} with the
\emph{arbitrary-portfolio} labor solutions (generalizing the
Proposition~\ref{prop:prices} labor block away from symmetric portfolios), and
then applies the method of undetermined coefficients at the symmetric-portfolio
point.

\subsection{Statement}

\begin{proposition}[Equilibrium portfolios]
\label{prop:portfolios}
Consider the simplified economy of Definition~\ref{def:simplified} with wages set
one period in advance and a positive steady-state labor wedge $\lwedge>0$ in each
country, and impose complete home bias $\hbias\to\zs/(1+\zs)$. Then, at least
around the symmetric-country-portfolio case:
\begin{enumerate}[label=(\alph*),leftmargin=1.8em,itemsep=3pt]
  \item Home's equilibrium portfolio share in capital (the world equity claim) is
  \emph{unaffected} by the supply of safe dollar debt, $\dfrac{dk}{db^g_{H,s}}=0$,
  while Home's portfolio share in dollar bonds \emph{falls} (Home issues more
  dollar bonds) as $-b^g_{H,s}$ rises, $\dfrac{db_H}{db^g_{H,s}}>0$;
  \item holding $\rra+\frac{1}{\zs}\rraF$ fixed, Home's portfolio share in capital
  \emph{rises} as Home becomes more risk-tolerant relative to Foreign,
  $\dfrac{dk}{d[\rraF/\rra]}\Big|_{\rra+\frac{1}{\zs}\rraF}>0$, while Home's
  portfolio share in dollar bonds \emph{falls} (Home issues more dollar bonds),
  $\dfrac{db_H}{d[\rraF/\rra]}\Big|_{\rra+\frac{1}{\zs}\rraF}<0$.
\end{enumerate}
\end{proposition}

\gapflag{We transcribe Proposition~4 to match KL's main-paper statement
(p.~1664) as relayed in the work-package brief; the project does not hold a local
copy of the main-paper PDF, so the verbatim wording is taken from the brief. The
mathematical content of (a)--(b) --- the four signed derivatives
$dk/db^g_{H,s}=0$, $db_H/db^g_{H,s}>0$, $dk/d[\rraF/\rra]>0$,
$db_H/d[\rraF/\rra]<0$ at the symmetric portfolio, the latter two holding
$\rra+\frac{1}{\zs}\rraF$ fixed --- matches the comparative statics displayed in
the appendix proof (their app.~p.~78), which we hold directly. Here $k$ denotes
Home's portfolio share in capital (the levered world-equity position) and $b_H$
Home's portfolio share in safe dollar bonds; ``Home issues more dollar bonds''
corresponds to $b_H$ falling / $-b^g_{H,s}$ rising.}

\subsection{Proof}

\paragraph{Strategy.} The proof has three steps. (i) Solve the labor block for
\emph{arbitrary} portfolios, generalizing the Proposition~\ref{prop:prices}
solution; this expresses $(\dev{\ell}_1,\fr{\dev{\ell}}_1)$ in terms of the shocks
$(\dev{\cy}_1,\dev{\epsilon}^z_1)$ and the forward state $\Eone\dev{\wsh}_2$. (ii)
Substitute into the international-portfolio equation~\eqref{eq:eq99} and into the
realized wealth-share identity~\eqref{eq:wsh-realized} to obtain a fixed-point in
$\Eone\dev{\wsh}_2$, yielding the eq.~(99)-with-$\Gamp$ identity. (iii) At the
symmetric-portfolio point (where a normalization constant $\Delta=1$), match
coefficients on $\dev{\epsilon}^z_1$ and $\dev{\cy}_1$ by the method of
undetermined coefficients to read off the portfolio derivatives.

\subparagraph{Step 1: arbitrary-portfolio labor solutions.} Generalizing the
Proposition~\ref{prop:prices} labor block (their app.~B.3.1) to arbitrary
portfolios, KL show (their app.~p.~77) that the symmetric and relative labor
movements solve
\begin{align}
\frac{1}{1+\zs}\dev{\ell}_1 + \frac{\zs}{1+\zs}\fr{\dev{\ell}}_1
&= -\frac{1}{\cshare+\frac{1}{\taylor}(1-\cshare)}\frac{1}{\taylor}\frac{1}{1+\zs}\dev{\cy}_1
+\frac{1}{\cshare+\frac{1}{\taylor}(1-\cshare)}(1-\cshare)\tel^z\dev{\epsilon}^z_1, \label{eq:lab-sym}\\
\fr{\dev{\ell}}_1 - \dev{\ell}_1
&= \frac{1}{\cshare+\frac{1}{\taylor}(1-\cshare)}\frac{1}{\taylor}\dev{\cy}_1
-\frac{1}{\cshare+\frac{1}{\taylor}(1-\cshare)}\frac{1+\zs}{\zs}\atil\,\Eone\dev{\wsh}_2. \label{eq:lab-rel}
\end{align}
These generalize the Proposition~1 labor solution: \eqref{eq:lab-sym} is the
aggregate-employment response (a positive safety shock contracts global
employment; a positive productivity surprise expands it), and \eqref{eq:lab-rel}
is the relative-employment response (a positive safety shock raises Foreign-less-
Home employment, since the shock hits the Home Taylor rule; and a positive Home
wealth share raises Home relative employment through the demand channel, with the
$\frac{1+\zs}{\zs}\atil$ wealth-share multiplier). The denominator
$\cshare+\frac{1}{\taylor}(1-\cshare)$ is the same $D$ of the
Proposition~\ref{prop:returns} proof.
\gapflag{KL display \eqref{eq:lab-sym}--\eqref{eq:lab-rel} as ``For arbitrary
portfolios, it is straightforward to show that [\dots], generalizing the results
given in the proof of Proposition~1 to arbitrary portfolios.'' We transcribe
their two displayed lines verbatim (app.~p.~77). The only difference from the
Proposition~1 labor block is the appearance of the $\Eone\dev{\wsh}_2$ term in
\eqref{eq:lab-rel}: under symmetric portfolios $\Eone\dev{\wsh}_2=0$ and
\eqref{eq:lab-rel} collapses to the Proposition~1 relative-labor response; for
arbitrary portfolios $\Eone\dev{\wsh}_2\neq0$ and the wealth share feeds back into
relative labor. We did not re-derive the labor block from the firm/wage system
\eqref{eq:firm-H}--\eqref{eq:firm-F}; we take KL's displayed arbitrary-portfolio
solution.}

\subparagraph{Step 2: substitute into the international-portfolio equation.}
Substituting \eqref{eq:lab-sym}--\eqref{eq:lab-rel} into the
international-portfolio equation~\eqref{eq:eq99}, and using the realized
wealth-share identity~\eqref{eq:wsh-realized} (their (92)) for $\Eone\dev{\wsh}_2$
on the left-hand side, KL obtain (their app.~pp.~76--77) the eq.~(99)-with-$\Gamp$
identity
\begin{multline}
\Bigl(\frac{\qk k}{\sav}-1\Bigr)\bigl(\rk_1-\rr_1\bigr)
+\frac{b_F}{\sav}\bigl(\fr{r}_1-\dev{\rer}_1-\rr_1\bigr)
-\disc\,\frac{\zs}{1+\zs}\,\frac{b^g_{H,s}}{\sav}\,\dev{\cy}_1
=\\
\frac{1}{\Gamp}(\rraF-\rra)(1-\cshare)\left[1+\lwedge(1-\disc)\frac{1}{\cshare+\frac{1}{\taylor}(1-\cshare)}(1-\cshare)\right]\tel^z\dev{\epsilon}^z_1\\
-\frac{1}{\Gamp}(\rraF-\rra)\lwedge(1-\disc)(1-\cshare)\frac{1}{\cshare+\frac{1}{\taylor}(1-\cshare)}\frac{1}{\taylor}\frac{1}{1+\zs}\dev{\cy}_1\\
+\frac{1}{\Gamp}\Bigl(\tfrac{1}{\zs}(\rraF-1)+(\rra-1)\Bigr)(1-\disc)\lwedge(1-\cshare)\frac{\zs}{1+\zs}\frac{1}{\cshare+\frac{1}{\taylor}(1-\cshare)}\frac{1}{\taylor}\dev{\cy}_1, \label{eq:eq99-prime}
\end{multline}
where the composite coefficient is now
\begin{equation}
\Gamp \;\equiv\; \Gam + \Bigl(\tfrac{1}{\zs}(\rraF-1)+(\rra-1)\Bigr)(1-\disc)\lwedge(1-\cshare)\frac{1}{\cshare+\frac{1}{\taylor}(1-\cshare)}\cshare \;>\;0. \label{eq:Gamma-prime-def}
\end{equation}
\gapflag{\eqref{eq:eq99-prime}--\eqref{eq:Gamma-prime-def} are KL's displayed
eq.~(99)-with-$\Gamp$ and the definition of $\Gamp$ (their app.~p.~77). The
reconstruction: substituting the relative-labor solution~\eqref{eq:lab-rel} into
the $(\fr{\dev{\ell}}_1-\dev{\ell}_1)$ terms of \eqref{eq:eq99} introduces a new
$\Eone\dev{\wsh}_2$ loading on the RHS (through the
$-\frac{1}{\cshare+\frac{1}{\taylor}(1-\cshare)}\frac{1+\zs}{\zs}\atil\Eone\dev{\wsh}_2$
piece of \eqref{eq:lab-rel}); moving that loading to the LHS augments the
$\Eone\dev{\wsh}_2$ coefficient $\Gam$ by exactly
$(\frac{1}{\zs}(\rraF-1)+(\rra-1))(1-\disc)\lwedge(1-\cshare)\frac{1}{\cshare+\frac{1}{\taylor}(1-\cshare)}\cshare$,
which is the added term in $\Gamp$. (The $\frac{1+\zs}{\zs}\atil$ from
\eqref{eq:lab-rel} cancels against the $\frac{\zs}{1+\zs}$ and one power of
$\atil$ in the relative-labor coefficient, leaving the $\cshare$ shown; the
positivity $\Gamp>0$ holds because, with $\rra,\rraF\ge1$ and $\lwedge\in(0,1)$,
the added term is non-negative and $\Gam>0$.) Substituting the symmetric-labor
solution~\eqref{eq:lab-sym} into the aggregate-employment term of \eqref{eq:eq99}
produces the $\dev{\epsilon}^z_1$ and $\dev{\cy}_1$ loadings shown on the RHS: the
productivity term collects the direct $(\rraF-\rra)(1-\cshare)$ loading plus the
wedge contribution $(\rraF-\rra)\lwedge(1-\disc)(1-\cshare)\times[\text{symmetric-labor coefficient on }\tel^z\dev{\epsilon}^z_1]$,
giving the displayed $[1+\lwedge(1-\disc)\frac{1-\cshare}{\cshare+\frac{1}{\taylor}(1-\cshare)}]$
bracket; the safety-shock terms collect the wedge and relative-labor
contributions through the $\dev{\cy}_1$ loadings of
\eqref{eq:lab-sym}--\eqref{eq:lab-rel}. We have matched KL's displayed eq.~(99)$'$
(app.~p.~77) line-by-line and flag the substitution as a genuine reconstruction:
KL display only ``Substituting these into the expression for international
portfolios [\dots] we obtain'' followed by the collected result. We reproduce the
collected result; the comparative statics below depend only on \emph{which} RHS
terms carry $\dev{\epsilon}^z_1$ versus $\dev{\cy}_1$ and on the sign of their
coefficients, which are robust to any residual algebra in the exact magnitudes.}

\subparagraph{Step 3: the symmetric-portfolio normalization and undetermined
coefficients.} The left-hand side of \eqref{eq:eq99-prime} is, by the realized
wealth-share identity~\eqref{eq:wsh-realized}, equal to $\Eone\dev{\wsh}_2$.
Solving \eqref{eq:eq99-prime} for $\Eone\dev{\wsh}_2$ and substituting the
realized returns of Proposition~\ref{prop:returns} (which, at the symmetric
portfolio, are the closed forms of Section~\ref{sec:engine}), KL express
$\Eone\dev{\wsh}_2$ as a linear combination of $\dev{\cy}_1$ and
$\dev{\epsilon}^z_1$ (their app.~p.~78):
\begin{multline}
\Eone\dev{\wsh}_2 =
\frac{1}{\Delta}\left[\frac{1}{\taylor}+\Bigl(1-\frac{1}{\taylor}\Bigr)\frac{\frac{1}{\taylor}(1-\cshare)}{\cshare+\frac{1}{\taylor}(1-\cshare)}\right]\frac{b_H}{\sav}\,\dev{\cy}_1
-\frac{1}{\Delta}\disc\,\frac{\zs}{1+\zs}\frac{b^g_{H,s}}{\sav}\,\dev{\cy}_1\\
+\frac{1}{\Delta}\Bigl(\frac{\qk k}{\sav}-1\Bigr)\left[(1-\cshare)+\Bigl(1-\frac{1}{\taylor}\Bigr)(1-\cshare)^2\frac{1}{\cshare+\frac{1}{\taylor}(1-\cshare)}\right]\tel^z\dev{\epsilon}^z_1, \label{eq:wsh-final}
\end{multline}
where
\begin{equation}
\Delta \;\equiv\; 1 + \left[\Bigl(\frac{\qk k}{\sav}-1\Bigr)+\frac{b_F}{\sav}\frac{\zs}{1+\zs}\right]
\frac{\frac{1}{\taylor}(1-\cshare)}{\cshare+\frac{1}{\taylor}(1-\cshare)}, \label{eq:Delta-def}
\end{equation}
which evaluates to $\Delta=1$ at the symmetric portfolio (where the bracketed
portfolio weights vanish). Here $b_H$ is Home's safe-dollar-bond position.
\gapflag{\eqref{eq:wsh-final}--\eqref{eq:Delta-def} are KL's displayed lines
(their app.~p.~78). The reconstruction: substituting the realized excess returns
$\rk_1-\rr_1$ and $\fr{r}_1-\dev{\rer}_1-\rr_1$ (Proposition~\ref{prop:returns},
both $\propto\dev{\cy}_1$ at the symmetric portfolio plus productivity loadings)
into the LHS of \eqref{eq:eq99-prime} makes the LHS itself a linear function of
$\Eone\dev{\wsh}_2$ (through the portfolio weights times the
$\Eone\dev{\wsh}_2$-dependence of the realized returns); collecting that
dependence onto the LHS produces the normalization factor $\Delta$ in
\eqref{eq:Delta-def}. At the symmetric portfolio $\frac{\qk k}{\sav}-1=0$ and
$\frac{b_F}{\sav}=0$, so $\Delta=1$ and the realized returns reduce to their
symmetric-portfolio forms. We transcribe KL's $\Delta$ and the
$\Eone\dev{\wsh}_2$ expression; the $\frac{b_H}{\sav}$ coefficient on the
$\dev{\cy}_1$ term comes from the leverage and Foreign-bond excess returns
(both $\propto\dev{\cy}_1$) collected through their portfolio weights, plus the
seigniorage term; the $(\frac{\qk k}{\sav}-1)$ coefficient on the
$\tel^z\dev{\epsilon}^z_1$ term comes from the productivity loading of
$\rk_1-\rr_1$. We flag that we reproduced KL's coefficients rather than
re-deriving each from the realized-return reduced forms.}

\medskip
We now read off the comparative statics by the method of undetermined
coefficients, evaluating at the symmetric portfolio ($\Delta=1$). Write the
equilibrium portfolio as the pair $(k,b_H)$ that makes \eqref{eq:wsh-final} hold
for \emph{both} independent shocks $\dev{\cy}_1$ and $\dev{\epsilon}^z_1$
simultaneously; equivalently, the coefficients on $\dev{\cy}_1$ and on
$\tel^z\dev{\epsilon}^z_1$ in \eqref{eq:wsh-final} must each equal the
corresponding coefficient required by the risk-sharing RHS of
\eqref{eq:eq99-prime}.

\emph{Comparative statics in $b^g_{H,s}$.} The supply of safe debt $b^g_{H,s}$
enters \eqref{eq:wsh-final} only through the seigniorage term
$-\frac{1}{\Delta}\disc\frac{\zs}{1+\zs}\frac{b^g_{H,s}}{\sav}\dev{\cy}_1$ and,
via the household budget, through $b_H$ (more government safe debt is held by
households). It does \emph{not} enter the $\tel^z\dev{\epsilon}^z_1$ coefficient,
which depends only on the capital weight $\frac{\qk k}{\sav}-1$. Matching the
$\dev{\epsilon}^z_1$ coefficient therefore leaves the equilibrium capital share
$k$ unchanged as $b^g_{H,s}$ varies:
\begin{equation}
\frac{dk}{db^g_{H,s}} = 0. \label{eq:cs-k-bg}
\end{equation}
Matching the $\dev{\cy}_1$ coefficient, an increase in $-b^g_{H,s}$ (more safe
debt issued) raises the seigniorage Home earns on impact of a positive safety
shock; to keep the risk-sharing condition satisfied at unchanged $k$, the
household must hold more dollar bonds, i.e.\ $b_H$ adjusts so that
\begin{equation}
\frac{db_H}{db^g_{H,s}} > 0. \label{eq:cs-bH-bg}
\end{equation}
\gapflag{Equations~\eqref{eq:cs-k-bg}--\eqref{eq:cs-bH-bg} are KL's displayed
comparative statics (app.~p.~78). The undetermined-coefficients argument:
\eqref{eq:wsh-final} must hold as an identity in the two independent shocks, so
the $\dev{\cy}_1$-coefficient equation and the $\dev{\epsilon}^z_1$-coefficient
equation are two scalar equations in the two unknown portfolio shares $(k,b_H)$.
Because $b^g_{H,s}$ appears \emph{only} in the $\dev{\cy}_1$ equation (the
seigniorage term), the $\dev{\epsilon}^z_1$ equation pins $k$ independently of
$b^g_{H,s}$, giving \eqref{eq:cs-k-bg}; differentiating the $\dev{\cy}_1$ equation
then gives the sign of $db_H/db^g_{H,s}$. We supply the
$2\times2$-block-triangular structure that makes $dk/db^g_{H,s}=0$ exact (the
$\dev{\epsilon}^z_1$ equation does not involve $b^g_{H,s}$); the sign of
$db_H/db^g_{H,s}$ follows from the positive seigniorage loading. KL state the two
results without displaying the linear-algebra block structure; we make it
explicit and flag this as a reconstruction of the undetermined-coefficients
step.}

\emph{Comparative statics in $\rraF/\rra$, holding $\rra+\frac{1}{\zs}\rraF$
fixed.} The combination $\rra+\frac{1}{\zs}\rraF$ enters $\Gam$ (and hence
$\Gamp$) in the denominator of every RHS term of \eqref{eq:eq99-prime}; holding it
fixed therefore freezes $\Gamp$. As KL note (their app.~p.~77), ``in comparative
statics with respect to $\rraF/\rra$ holding fixed $\rra+\frac{1}{\zs}\rraF$, on
the right-hand side it is clear that only the first two terms vary; the third,
capturing the effect of the relative labor response on international portfolios,
is constant.'' The first two RHS terms both carry the factor $(\rraF-\rra)$: the
productivity term $\frac{1}{\Gamp}(\rraF-\rra)(1-\cshare)[\cdots]\tel^z\dev{\epsilon}^z_1$
and the aggregate-employment safety term
$-\frac{1}{\Gamp}(\rraF-\rra)\lwedge(1-\disc)(1-\cshare)[\cdots]\dev{\cy}_1$. An
increase in $\rraF/\rra$ at fixed $\rra+\frac{1}{\zs}\rraF$ raises $\rraF-\rra$.
Matching the $\dev{\epsilon}^z_1$ coefficient, the productivity loading on the RHS
rises, so the equilibrium capital weight must rise to deliver a larger Home
wealth gain in good productivity states:
\begin{equation}
\frac{dk}{d[\rraF/\rra]}\Big|_{\rra+\frac{1}{\zs}\rraF} > 0. \label{eq:cs-k-gamma}
\end{equation}
Matching the $\dev{\cy}_1$ coefficient (with $\lwedge>0$), the safety-shock
loading on the RHS becomes more negative, so to satisfy risk sharing Home must
issue more dollar bonds (its dollar-bond share falls):
\begin{equation}
\frac{db_H}{d[\rraF/\rra]}\Big|_{\rra+\frac{1}{\zs}\rraF} < 0
\qquad(\text{assuming } \lwedge>0). \label{eq:cs-bH-gamma}
\end{equation}
\gapflag{Equations~\eqref{eq:cs-k-gamma}--\eqref{eq:cs-bH-gamma} are KL's
displayed comparative statics (app.~p.~78), with the second holding
$\rra+\frac{1}{\zs}\rraF$ fixed and assuming $\lwedge>0$. The
undetermined-coefficients argument mirrors the $b^g_{H,s}$ case: the
$\dev{\epsilon}^z_1$ equation pins $k$ and, because its RHS coefficient rises with
$\rraF-\rra$ (the productivity term), $dk/d[\rraF/\rra]>0$; the $\dev{\cy}_1$
equation then pins $b_H$, and because its RHS coefficient (the
aggregate-employment safety term, which carries $-(\rraF-\rra)\lwedge$) becomes
more negative with $\rraF-\rra$ when $\lwedge>0$, $db_H/d[\rraF/\rra]<0$. We
supply the matching argument and flag (i) the reliance on KL's own statement that
only the first two RHS terms vary at fixed $\rra+\frac{1}{\zs}\rraF$, and (ii) the
$\lwedge>0$ requirement for the sign of $db_H/d[\rraF/\rra]$, which KL state
explicitly (``assumes $\lwedge>0$''). The signs depend on the positivity of
$\Gamp$ \eqref{eq:Gamma-prime-def} and on the steady-state portfolio being levered
long capital and short dollar bonds, as at the symmetric benchmark.}

\medskip
This establishes parts (a) and (b) of Proposition~\ref{prop:portfolios}. \qed
(\,Proposition~\ref{prop:portfolios}\,)

\subsection{Discussion}

\paragraph{Two risk factors, two portfolio margins.} Proposition~4 organizes the
equilibrium portfolio around the two shocks the asset span must insure: a
\emph{global productivity} factor $\dev{\epsilon}^z_1$ and a \emph{safe-asset
demand} (safety) factor $\dev{\cy}_1$. The two portfolio margins --- the capital
(world-equity) share $k$ and the dollar-bond share $b_H$ --- map to the two
factors through the block structure exposed in the proof: the capital share is
pinned by the productivity-coefficient equation, the dollar-bond share by the
safety-coefficient equation.

\paragraph{The role of safe-debt supply.} As Home's government borrows more in
safe dollar bonds ($-b^g_{H,s}$ rises), Home receives more seigniorage on impact
of a positive safety shock: it is the issuer of the asset whose convenience value
spikes exactly when safety is demanded. This makes Home a \emph{natural insurer}
of the safety shock, and --- crucially --- it can provide that insurance
\emph{without loading up on productivity risk}. Concretely, Home borrows more in
dollar bonds (lower $b_H$) to hold more Foreign bonds, leaving its capital
(world-equity) exposure $k$ unchanged: this is exactly the content of
$dk/db^g_{H,s}=0$ together with $db_H/db^g_{H,s}>0$. The supply of safe dollar
debt is a margin for insuring the safety factor that is orthogonal to the
equity-risk margin.

\paragraph{The role of relative risk tolerance.} As Home becomes more
risk-tolerant than Foreign ($\rraF/\rra$ rises, holding
$\rra+\frac{1}{\zs}\rraF$ fixed so that the aggregate price of risk $\Gam$ is
held constant), Home optimally insures \emph{both} the negative-productivity and
the positive-safety shocks: it holds more capital (so it bears more of the global
equity risk, $dk/d[\rraF/\rra]>0$) and it borrows more in dollar bonds (so it
provides more safety insurance, $db_H/d[\rraF/\rra]<0$). KL emphasize that this is
the sense in which \emph{greater risk tolerance is necessary} to explain why the
United States takes a disproportionate exposure to equity returns: in their model,
the US equity-tilted external balance sheet requires the US to be more
risk-tolerant than the rest of the world.

\paragraph{The big\_cy reading.} Proposition~4 is precisely the KL tension that
big\_cy aims to revisit: KL need greater US risk tolerance ($\rra<\rraF$) to
rationalize the US equity tilt, yet high US risk tolerance is the very feature the
project regards as counterfactual. The two margins isolated here ---
$dk/d[\rraF/\rra]>0$ (equity tilt needs risk tolerance) and $db_H/db^g_{H,s}>0$
(safe-debt supply moves the dollar-bond margin without touching equity) ---
identify exactly where a large \emph{absolute} convenience yield might substitute
for the risk-tolerance asymmetry. We record KL's interpretation faithfully here
and defer the big\_cy extension to the dedicated extension memo.
